\subsection{고정 게임 native runtime 비교}
이 공개 결과는 하나의 authored quest(\texttt{sealed-lighthouse-v1})에 대해 seed 20261024, selector당 representation별 반복 2회의 8개 배정 행을 보존한다. 시스템 endpoint 성공은 5/8이고, 관측된 native 실행은 7개이다. 실제 receipt에서 집계한 host 호출은 23회이며, 관측된 repair 합은 0이다. committed transition은 16, 그중 unsafe는 0이며, semantic replay valid/invalid/NR은 7/0/1이다. fresh-engine reexecution 처분은 VALID=6; INVALID=0; NOT\_RUN=1; NOT\_OBSERVED=1이다. elapsed 관측은 n=7; 48.19/41.09초(평균/중앙값), preparation 관측은 n=7; 6.67/6.30초이다. 보고되지 않은 값은 0이 아니라 NR이다.

첫 행은 harness JSON serialization 기술 실패로 유지되어 GPT-6.1 능력에 귀속하지 않는다. GPT-6.1 ontology에서 response-timeout 행 2개 이상을 frozen 60초 경계 안의 실패로 보존한다. 모델은 typed-intent proposal만 담당하고 deterministic Python planner/checker와 Godot commit이 권한을 유지한다. 따라서 이 표는 한 quest의 기술적 관측이며 efficacy, superiority, noninferiority, 일반 ontology 일관성 또는 player benefit을 주장하지 않는다. semantic record-link replay만으로 fresh-engine reexecution을 주장하지 않는다. provider identity/revision/cost와 provider request 수는 공개 관측 경계를 넘어 NR이다.
